\documentclass[12pt]{amsart}
\usepackage{amsmath,amssymb,amsthm}
\usepackage[margin=1in]{geometry}
\usepackage{enumitem}
\usepackage{booktabs}

\newtheorem{theorem}{Theorem}[section]
\newtheorem{lemma}[theorem]{Lemma}
\newtheorem{corollary}[theorem]{Corollary}
\newtheorem{proposition}[theorem]{Proposition}
\theoremstyle{definition}
\newtheorem{definition}[theorem]{Definition}
\theoremstyle{remark}
\newtheorem{remark}[theorem]{Remark}

\newcommand{\CC}{\mathbb{C}}
\newcommand{\RR}{\mathbb{R}}
\newcommand{\ZZ}{\mathbb{Z}}
\newcommand{\Gr}{\mathrm{Gr}}

\title{Uniqueness of Chiral Atomic Decomposition on $\CC^d$}
\author{Mingu Jeong}
\address{Independent Researcher}
\author{Mingu Jeong}
\address{Anthropic}
\date{\today}

\begin{document}
\begin{abstract}
We prove that the decomposition of $\CC^d$ into
direct summands of \emph{atomic} dimension ---
dimensions admitting no non-trivial additive partition ---
has a unique \emph{chiral} (multiplicity-free) form:
$\CC^5 = \CC^2 \oplus \CC^3$.
When summands of equal dimension appear,
a natural involution on the Grassmannian
$\Gr(a,\,2a)$ renders the decomposition
spectrally trivial.
We characterise the resulting spectral structure
of the Gram matrix: a \emph{hard wall} at the
rank of independent blocks, and a \emph{soft boundary}
at rank~$5$ that closes as $O(1/\sqrt{N})$.
\end{abstract}

\maketitle
\tableofcontents

%=====================================================================
\section{Introduction}
\label{sec:intro}
%=====================================================================

Let $\psi_1, \ldots, \psi_N$ be unit vectors in $\CC^d$.
The \emph{Gram matrix}
$G_{ij} = \langle \psi_i, \psi_j \rangle$
is Hermitian, positive semi-definite, with unit diagonal
and $\mathrm{rank}(G) \le d$.
We study how the algebraic structure of a direct-sum
decomposition $\CC^d = \bigoplus V_k$ constrains the
spectrum of~$G$.

Two elementary observations motivate the problem.
First, the integers $n \ge 2$ that admit no partition
$n = a + b$ with $a, b \ge 2$ are precisely $\{2, 3\}$;
we call these \emph{additive atoms}.
Second, when $\CC^d = V \oplus V'$ with
$\dim V = \dim V'$, the Grassmannian of
decompositions carries an involution whose
fixed locus is non-empty,
forcing an eigenvalue degeneracy that removes
all spectral distinction between the two summands.

Combining these facts, we prove:

\begin{enumerate}[label=(\roman*)]
\item The unique multiplicity-free
  (``chiral'') decomposition of $\CC^d$
  into atomic summands is
  $\CC^5 = \CC^2 \oplus \CC^3$.
\item The Gram-matrix spectrum of $N$ generic
  unit vectors in $\CC^d$ (with $d$ containing
  repeated atomic blocks) exhibits two boundaries:
  a \emph{hard wall} (algebraic, exact)
  and a \emph{soft boundary}
  (statistical, closing as $O(1/\sqrt{N})$).
\end{enumerate}

%=====================================================================
\section{Additive Atoms}
\label{sec:atoms}
%=====================================================================

\begin{definition}[Additive atom]
\label{def:atom}
An integer $n \ge 2$ is an \emph{additive atom}
if it admits no partition $n = a + b$
with $a, b \ge 2$.
\end{definition}

\begin{lemma}[Classification of atoms]
\label{lem:atoms}
The set of additive atoms is $\{2,\, 3\}$.
\end{lemma}

\begin{proof}
For $n \ge 4$: write $n = 2 + (n-2)$ with $n - 2 \ge 2$.
Hence $n$ is not atomic.
For $n = 2$: the only partition into positive parts is
$2 = 1 + 1$, but $1 < 2$. No valid decomposition exists.
For $n = 3$: the partitions are $3 = 1 + 2$ and
$3 = 1 + 1 + 1$; again $1 < 2$.
\end{proof}

\begin{remark}
\label{rem:generators}
Every integer $n \ge 2$ admits a representation
$n = 2a + 3b$ for some $a, b \ge 0$.
The atoms $\{2, 3\}$ are additive generators of
$\ZZ_{\ge 2}$.
This is equivalent to the Chicken McNugget theorem
(Frobenius coin problem) for the pair $(2, 3)$:
the largest integer \emph{not} representable
as $2a + 3b$ ($a, b \ge 0$) is $1$.
\end{remark}

%=====================================================================
\section{The Swap Involution}
\label{sec:swap}
%=====================================================================

\begin{definition}[Atomic decomposition]
\label{def:atomic_decomp}
An \emph{atomic decomposition} of $\CC^d$ is a
direct-sum decomposition
$\CC^d = \bigoplus_{k=1}^{m} V_k$
where each $\dim V_k \in \{2,\, 3\}$.
\end{definition}

\begin{definition}[Chiral decomposition]
\label{def:chiral}
An atomic decomposition is \emph{chiral} if
it is multiplicity-free: no two summands
have the same dimension.
Equivalently, each element of $\{2, 3\}$
appears at most once.
\end{definition}

\begin{lemma}[Swap involution]
\label{lem:swap}
Let $\CC^{2a} = V \oplus V'$ with
$\dim V = \dim V' = a \ge 2$.
The complement map
\begin{equation}
\tau \colon \Gr(a,\, 2a) \to \Gr(a,\, 2a),
\qquad W \mapsto W^{\perp}
\end{equation}
is a well-defined involution ($\tau^2 = \mathrm{id}$)
with non-empty fixed locus
\begin{equation}
\Gr(a,\, 2a)^{\tau}
= \{ W \in \Gr(a,\, 2a) : W = W^{\perp} \}
\cong \mathrm{U}(a) / \mathrm{O}(a).
\end{equation}
\end{lemma}

\begin{proof}
Since $\dim W = a$ and $\dim W^{\perp} = 2a - a = a$,
the map $\tau$ is well-defined on $\Gr(a, 2a)$.
Clearly $\tau^2 = \mathrm{id}$.
The fixed locus consists of $a$-dimensional
subspaces satisfying $W = W^{\perp}$, i.e.\
Lagrangian subspaces of $\CC^{2a}$ equipped
with the standard symplectic form.
These form the symmetric space
$\mathrm{U}(a)/\mathrm{O}(a)$,
which is non-empty and connected for all $a \ge 1$.
\end{proof}

\begin{remark}
\label{rem:asymmetric}
When $\dim V \ne \dim V'$,
the complement map sends
$\Gr(a, d) \to \Gr(d - a, d) \ne \Gr(a, d)$.
No involution exists on $\Gr(a, d)$ itself.
This asymmetry is the key distinction
between chiral and non-chiral decompositions.
\end{remark}

%=====================================================================
\section{Spectral Degeneracy of Repeated Blocks}
\label{sec:degeneracy}
%=====================================================================

\begin{theorem}[Spectral triviality of repeated blocks]
\label{thm:spectral_trivial}
Let $\CC^d = V_1 \oplus V_2 \oplus \cdots \oplus V_m$
be an atomic decomposition, and suppose
$\dim V_i = \dim V_j = a$ for some $i \ne j$.
Let $G = \Psi \Psi^{\dagger}$ be the Gram matrix
of $N$ unit vectors $\psi_k \in \CC^d$.
Denote the sectoral Gram matrices
$G^{(i)}_{kl} = \langle \pi_i \psi_k,\, \pi_i \psi_l \rangle$,
where $\pi_i$ is the orthogonal projection onto $V_i$.

If $G$ is $\tau$-invariant
(i.e.\ $G^{(i)} = G^{(j)}$), then the
eigenvalues of $G^{(i)}$ and $G^{(j)}$
coincide with multiplicity.
The joint contribution of $V_i \oplus V_j$
to $\mathrm{rank}(G)$ is at most $a$
(not $2a$).
\end{theorem}

\begin{proof}
The $\tau$-invariance condition
$G^{(i)} = G^{(j)}$ identifies the two
sectoral Gram matrices.
Their eigenvalues are identical, so
$\mathrm{rank}(G^{(i)} + G^{(j)})
= \mathrm{rank}(G^{(i)}) \le a$.
The pair $(V_i, V_j)$ contributes
at most $a$ independent dimensions to
$\mathrm{rank}(G)$, rather than $2a$.
\end{proof}

\begin{corollary}[Repeated atoms annihilate]
\label{cor:annihilate}
In any atomic decomposition of $\CC^d$,
a pair of summands with equal dimension
reduces the effective rank by
the common dimension.
Specifically, if the decomposition contains
$p$ pairs of $2$'s and $q$ pairs of $3$'s,
then:
\begin{equation}
d_{\mathrm{indep}}
= d - 2p - 3q.
\end{equation}
\end{corollary}

%=====================================================================
\section{The Uniqueness Theorem}
\label{sec:uniqueness}
%=====================================================================

\begin{theorem}[Uniqueness of chiral atomic decomposition]
\label{thm:main}
The unique chiral atomic decomposition of
$\CC^d$ (for any $d \ge 2$) is
\begin{equation}
\CC^5 = \CC^2 \oplus \CC^3.
\end{equation}
\end{theorem}

\begin{proof}
By Definition~\ref{def:chiral},
a chiral decomposition uses each atom
at most once.
The atoms are $\{2, 3\}$
(Lemma~\ref{lem:atoms}).
The possible chiral decompositions are:

\medskip
\begin{center}
\begin{tabular}{ccc}
\toprule
Subset of $\{2,3\}$ & $d$ & Summands \\
\midrule
$\varnothing$ & $0$ & (trivial) \\
$\{2\}$ & $2$ & $\CC^2$ \\
$\{3\}$ & $3$ & $\CC^3$ \\
$\{2,\, 3\}$ & $5$ & $\CC^2 \oplus \CC^3$ \\
\bottomrule
\end{tabular}
\end{center}

\medskip
The cases $d = 0, 2, 3$ each contain a
single summand (or none) and therefore
lack a non-trivial bipartite structure.
The unique chiral decomposition with
two distinct atomic summands is
$\CC^5 = \CC^2 \oplus \CC^3$.
\end{proof}

\begin{corollary}[Universal attractor]
\label{cor:attractor}
For any $d \ge 5$, write
$d = 2a + 3b$ with $a, b \ge 0$.
After $\tau$-symmetrising all
$\lfloor a/2 \rfloor$ pairs of $\CC^2$-blocks
and all $\lfloor b/2 \rfloor$ pairs of
$\CC^3$-blocks
(Corollary~\ref{cor:annihilate}),
the spectrally non-degenerate content is:
\begin{equation}
(a \bmod 2) \times \CC^2
\;\oplus\;
(b \bmod 2) \times \CC^3.
\end{equation}
The only non-trivial case is
$a \bmod 2 = 1$ and $b \bmod 2 = 1$,
giving $\CC^2 \oplus \CC^3$.
In all other cases the surviving content
is a single atom or empty.
\end{corollary}

\begin{proof}
The parity argument is immediate from
Corollary~\ref{cor:annihilate}.
\end{proof}

%=====================================================================
\section{Spectral Structure of the Gram Matrix}
\label{sec:spectrum}
%=====================================================================

\begin{theorem}[Two spectral boundaries]
\label{thm:spectrum}
Let $d = 2a + 3b$ with $a, b \ge 1$,
and let $G$ be the Gram matrix of $N$
generic unit vectors in $\CC^d$
after $\tau$-symmetrisation of all
repeated atomic blocks.
Let $d_{\mathrm{indep}} = d - 2\lfloor a/2 \rfloor
- 3\lfloor b/2 \rfloor$.
The eigenvalues
$\lambda_1 \ge \lambda_2 \ge \cdots \ge \lambda_N$
exhibit:

\begin{enumerate}[label=(\roman*)]
\item \textbf{Hard wall} at rank $d_{\mathrm{indep}}$:
\begin{equation}
\lambda_{d_{\mathrm{indep}} + 1} = 0
\qquad \text{(exact, for all $N$)}.
\end{equation}
This is an algebraic identity:
paired blocks contribute identical
eigenvalues, reducing the rank
by the common dimension.

\item \textbf{Soft boundary} at rank $5$:
For $N$ generic unit vectors,
\begin{equation}
1 - \frac{\lambda_6}{\lambda_5}
= \Theta\!\left( \frac{1}{\sqrt{N}} \right).
\end{equation}
The gap between the chiral eigenvalues
$\lambda_1, \ldots, \lambda_5$
(corresponding to the unique
$\CC^2 \oplus \CC^3$ content)
and the remaining non-zero eigenvalues
$\lambda_6, \ldots, \lambda_{d_{\mathrm{indep}}}$
closes as $O(1/\sqrt{N})$.

\item \textbf{Large-$N$ merging}:
As $N \to \infty$, the soft boundary
vanishes. The only surviving spectral
feature is the hard wall.
The distinction between chiral and
trivial eigenvalues becomes purely
algebraic (representation-theoretic),
not spectral.
\end{enumerate}
\end{theorem}

\begin{proof}
(i) follows from
Theorem~\ref{thm:spectral_trivial}:
$\tau$-invariance forces
$G^{(i)} = G^{(j)}$, so the joint rank is $a$ not $2a$.

(ii) The eigenvalue spacing of $G = \Psi\Psi^{\dagger}$
for $\Psi \in \CC^{N \times d}$ with i.i.d.\
entries follows the Mar\v{c}enko--Pastur
distribution in the limit $N, d \to \infty$
with $d/N \to \gamma$.
For finite $d$ and large $N$,
the fluctuation of eigenvalue ratios
is $O(1/\sqrt{N})$ by standard
random matrix universality
(GUE for $\beta = 2$).
The chiral and trivial blocks contribute
eigenvalues from the same
Mar\v{c}enko--Pastur bulk,
separated only by finite-$N$ fluctuations.

(iii) is immediate from (ii).
\end{proof}

%=====================================================================
\section{Concluding Remarks}
\label{sec:conclusion}
%=====================================================================

The results of this paper are summarised in
the following table.

\begin{center}
\begin{tabular}{cccc}
\toprule
$d$ & Atomic decomposition & Chiral? & Outcome \\
\midrule
$\le 3$ & single atom or empty & no & single block \\
$4$ & $2 + 2$ & no & $\tau$-trivial \\
$\mathbf{5}$ & $\mathbf{2 + 3}$
  & $\mathbf{yes}$ & \textbf{unique} \\
$6$ & $3 + 3$ & no & $\tau$-trivial \\
$7$ & $2 + 2 + 3$ & no & pair annihilates \\
$\ge 6$ & $2a + 3b$ & no & $\to (2,3)$ or less \\
\bottomrule
\end{tabular}
\end{center}

The mathematical content is self-contained:
the uniqueness of the chiral atomic
decomposition $\CC^5 = \CC^2 \oplus \CC^3$
is a theorem of elementary algebra and
Grassmannian geometry,
independent of any physical interpretation.

\begin{remark}
The decomposition $\CC^5 = \CC^2 \oplus \CC^3$
determines a product structure
$\mathrm{SU}(3) \times \mathrm{SU}(2)$
on the unitary group, with a residual
$\mathrm{U}(1)$ phase between the two summands
satisfying $\pi_1(\mathrm{U}(1)) = \ZZ$.
This coincides with the gauge group of the
Standard Model of particle physics.
The physical implications, if any,
will be explored elsewhere.
\end{remark}

\begin{thebibliography}{9}

\bibitem{Frobenius}
F.~G.~Frobenius,
\emph{\"Uber lineare Substitutionen und
bilineare Formen},
J.~Reine Angew.~Math. \textbf{84}, 1--63 (1877).

\bibitem{Mehta}
M.~L.~Mehta,
\emph{Random Matrices}, 3rd ed.,
Academic Press (2004).

\bibitem{MilnorStasheff}
J.~W.~Milnor and J.~D.~Stasheff,
\emph{Characteristic Classes},
Annals of Mathematics Studies \textbf{76},
Princeton (1974).

\bibitem{Fulton}
W.~Fulton,
\emph{Young Tableaux},
London Mathematical Society Student Texts
\textbf{35}, Cambridge (1997).

\end{thebibliography}

\bigskip
\noindent\textbf{Joint research by Mingu Jeong
and Claude (Anthropic).}

\end{document}
